% GENERATED FILE. Edit canonical lessons, then run npm run generate:books.
% canonical-source-hash: fnv1a64:9737baf1e3b12c24
% canonical-lessons: GE-C153-wandern, GE-C153-wechseln, GE-C153-werfen, GE-C153-wiederholen, GE-C153-ziehen

\chapter{To hike, To change, To throw, To repeat, To pull}
\label{ch:ge-to-hike-to-change-to-throw-to-repeat-to-pull}
\hlchaptermodality{153}

\begin{chapteropening}
\textbf{By the end of this chapter:} \emph{I can say to hike, to change, to throw, to repeat and to pull.}

Complete the last lesson of chapter 153: I can say to hike, to change, to throw, to repeat and to pull.
\end{chapteropening}

\section[wandern]{wandern — to hike}

\label{lesson:GE-C153-wandern}



\begin{quote}
Before the new one: say the German for to promise, then the German for to introduce.
\end{quote}

\subsection*{You'll want to know: wandern}
\textbf{wandern} — \textquotedblleft{}to hike\textquotedblright{}.

\begin{grammarlens}[title={What you've built}]
One more word for this chapter.
\end{grammarlens}

\subsection*{Your turn}
\begin{itemize}
\raggedright
  \item \emph{Say it:} \emph{wandern}
  \item \emph{Say it:} \emph{wandern}, once more
  \item \emph{Say it:} say \emph{wandern}
  \item \emph{Recall:} say the German for to promise, then the German for to introduce, then say \emph{wandern} again
\end{itemize}

\begin{tcolorbox}[breakable,colback=teal!4,colframe=teal!35!black,title={Before you move on}]
What does wandern mean? (\textquotedblleft{}To hike\textquotedblright{}.) Say it once more.
\end{tcolorbox}
\section[wechseln]{wechseln — to change (money), to switch}

\label{lesson:GE-C153-wechseln}



\begin{quote}
Before the new one: say the German for to introduce, then the German for to hike.
\end{quote}

\subsection*{You'll want to know: wechseln}
\textbf{wechseln} — \textquotedblleft{}to change (money), to switch\textquotedblright{}.

\begin{grammarlens}[title={What you've built}]
One more word for this chapter.
\end{grammarlens}

\subsection*{Your turn}
\begin{itemize}
\raggedright
  \item \emph{Say it:} \emph{wechseln}
  \item \emph{Say it:} \emph{wechseln}, once more
  \item \emph{Say it:} say \emph{wechseln}
  \item \emph{Recall:} say the German for to introduce, then the German for to hike, then say \emph{wechseln} again
\end{itemize}

\begin{tcolorbox}[breakable,colback=teal!4,colframe=teal!35!black,title={Before you move on}]
What does wechseln mean? (\textquotedblleft{}To change (money), to switch\textquotedblright{}.) Say it once more.
\end{tcolorbox}
\section[werfen]{werfen — to throw}

\label{lesson:GE-C153-werfen}



\begin{quote}
Before the new one: say the German for to hike, then the German for to change.
\end{quote}

\subsection*{You'll want to know: werfen}
\textbf{werfen} — \textquotedblleft{}to throw\textquotedblright{}.

\begin{grammarlens}[title={What you've built}]
One more word for this chapter.
\end{grammarlens}

\subsection*{Your turn}
\begin{itemize}
\raggedright
  \item \emph{Say it:} \emph{werfen}
  \item \emph{Say it:} \emph{werfen}, once more
  \item \emph{Say it:} say \emph{werfen}
  \item \emph{Recall:} say the German for to hike, then the German for to change, then say \emph{werfen} again
\end{itemize}

\begin{tcolorbox}[breakable,colback=teal!4,colframe=teal!35!black,title={Before you move on}]
What does werfen mean? (\textquotedblleft{}To throw\textquotedblright{}.) Say it once more.
\end{tcolorbox}
\section[wiederholen]{wiederholen — to repeat}

\label{lesson:GE-C153-wiederholen}



\begin{quote}
Before the new one: say the German for to change, then the German for to throw.
\end{quote}

\subsection*{You'll want to know: wiederholen}
\textbf{wiederholen} — \textquotedblleft{}to repeat\textquotedblright{}.

\begin{grammarlens}[title={What you've built}]
One more word for this chapter.
\end{grammarlens}

\subsection*{Your turn}
\begin{itemize}
\raggedright
  \item \emph{Say it:} \emph{wiederholen}
  \item \emph{Say it:} \emph{wiederholen}, once more
  \item \emph{Say it:} say \emph{wiederholen}
  \item \emph{Recall:} say the German for to change, then the German for to throw, then say \emph{wiederholen} again
\end{itemize}

\begin{tcolorbox}[breakable,colback=teal!4,colframe=teal!35!black,title={Before you move on}]
What does wiederholen mean? (\textquotedblleft{}To repeat\textquotedblright{}.) Say it once more.
\end{tcolorbox}
\section[ziehen]{ziehen — to pull}

\label{lesson:GE-C153-ziehen}



\begin{quote}
Before the new one: say the German for to throw, then the German for to repeat.
\end{quote}

\subsection*{You'll want to know: ziehen}
\textbf{ziehen} — \textquotedblleft{}to pull\textquotedblright{}.

\begin{grammarlens}[title={What you've built}]
One more word for this chapter.
\end{grammarlens}

\subsection*{Your turn}
\begin{itemize}
\raggedright
  \item \emph{Say it:} \emph{ziehen}
  \item \emph{Say it:} \emph{ziehen}, once more
  \item \emph{Say it:} say \emph{ziehen}
  \item \emph{Recall:} say the German for to throw, then the German for to repeat, then say \emph{ziehen} again
\end{itemize}

\begin{tcolorbox}[breakable,colback=teal!4,colframe=teal!35!black,title={Before you move on}]
What does ziehen mean? (\textquotedblleft{}To pull\textquotedblright{}.) Say it once more.
\end{tcolorbox}
